\section*{الفصل \ref{ch:g10:heart-lungs-effort} --- القلب والرئتان أثناء الجهد}

\begin{solution}{exo:g10:heart-lungs-effort:1}
$0.6 \times 14 = \qty{8.4}{L/min}$؛ و$2.2 \times 35 = \qty{77}{L/min}$.
\end{solution}

\begin{solution}{exo:g10:heart-lungs-effort:2}
حجم الدم الذي يضخه أحد جانبي القلب في الدقيقة:
$Q = f \times V_s$. و$65 \times 0.075 \approx \qty{4.9}{L/min}$.
\end{solution}

\begin{solution}{exo:g10:heart-lungs-effort:3}
الأذين الأيمن، فالبطين الأيمن، فالشريان الرئوي، فالرئتان (الدورة
الرئوية)، فالأوردة الرئوية، فالأذين الأيسر، فالبطين الأيسر، فالأبهر،
فالأعضاء (الدورة العامة)، فالأوردة، فالأذين الأيمن.
\end{solution}

\begin{solution}{exo:g10:heart-lungs-effort:4}
الحجم الجاري \qty{0.5}{L}؛ والسعة الحيوية \qty{4.5}{L}؛ والحجم
المتبقي \qty{1.2}{L}.
\end{solution}

\begin{solution}{exo:g10:heart-lungs-effort:5}
نحو $220 - 16 = 204$ و$220 - 60 = 160$ نبضة في الدقيقة.
\end{solution}

\begin{solution}{exo:g10:heart-lungs-effort:6}
عند \qty{100}{W}: $120 \times 0.105 \approx \qty{12.6}{L/min}$؛ وعند
\qty{300}{W}: $198 \times 0.115 \approx \qty{22.8}{L/min}$؛ أي بعامل
قدره 1.8.
\end{solution}

\begin{solution}{exo:g10:heart-lungs-effort:7}
الراحة: \qty{2.5}{L/min}؛ والجهد: $0.04 \times 24 \approx
\qty{1.0}{L/min}$. فالنصيب يهبط اثني عشر ضعفًا، أما الجريان المطلق
فبعامل 2.6 فقط، لأن الصبيب نفسه قد نما.
\end{solution}

\begin{solution}{exo:g10:heart-lungs-effort:8}
$18 \times 0.14 = \qty{2.5}{L/min}$: أي نحو ثلاثة أرباع
\qty{3.45}{L/min} عند تلميذ مدرَّب --- وهو قريب منها عند شخص أصغر.
\end{solution}

\begin{solution}{exo:g10:heart-lungs-effort:9}
الدماغ لا يحتمل حتى ثواني بلا أكسجين وغلوكوز، فإمداده محمي. أما المعي
والكليتان فتستطيعان تعليق الهضم والرشح ساعة بلا ضرر، ويعار نصيبهما
للعضلات.
\end{solution}

\begin{solution}{exo:g10:heart-lungs-effort:10}
$V_s = 5000/42 \approx \qty{120}{mL}$، مقابل نحو \qty{70}{mL}: فقلبها
المتضخم يوصل صبيب الراحة بنبضات أقل بكثير.
\end{solution}

\begin{solution}{exo:g10:heart-lungs-effort:11}
أمر الدماغ بالحركة، واستباقه للجهد، يرسلان إلى المراكز القلبية
والغدتين الكظريتين قبل أن تعمل العضلات. أما عدّاء لا يعتمد إلا على
تحسس ثاني أكسيد الكربون فكان ليبدأ بصبيب الراحة، ويجري الدقيقة الأولى
على \omterm{def:g10:cell-metabolism:fermentation}{التخمر}، ويضطر إلى الإبطاء بالحمض اللبني قبل أن يلحق به الإيصال.
\end{solution}

\begin{solution}{exo:g10:heart-lungs-effort:12}
$195 \times 0.110 \times 0.150 \approx \qty{3.2}{L/min}$ و
$195 \times 0.160 \times 0.150 \approx \qty{4.7}{L/min}$. والحجم
النفضي هو ما يكبّره التدريب: فالشخص الثاني هو المدرَّب على التحمل.
\end{solution}

\begin{solution}{exo:g10:heart-lungs-effort:13}
الأكسجين المستنشق: $0.21 \times 100 = \qty{21}{L/min}$؛ والمستعمل:
\qty{3.5}{L/min}، أي 17\%. أما الباقي فيزفر من جديد --- فالهواء
الزفيري لا يزال يحمل نحو 16\% أكسجينًا، ولذلك ينفع الإنعاش بالنفخ من
الفم إلى الفم.
\end{solution}

\begin{solution}{exo:g10:heart-lungs-effort:14}
الصبيب على الأكثر نحو $110 \times 0.12 \approx \qty{13}{L/min}$ بدل
24: فيهبط $\dot V\!\mathrm{O_2}$max قرابة النصف، إلى نحو
\qty{2}{L/min}. ويستطيع المريض المشي وركوب الدراجة برفق لكن لا إدامة
أي جهد شاق.
\end{solution}

\begin{solution}{exo:g10:heart-lungs-effort:15}
الإيصال في كل لتر يهبط الربع، والاستخلاص لا يتغير نسبيًا:
فيهبط $\dot V\!\mathrm{O_2}$max بنسبة 25\%. وفي الأسابيع التالية يصنع
الجسم كريات حمراء أكثر، فيرفع الأكسجين المحمول في اللتر عائدًا نحو
\qty{200}{mL}.
\end{solution}

\begin{solution}{pb:g10:heart-lungs-effort:1}
\textbf{1.} $62 \times 0.072 \approx 4.5$؛ و$105 \times 0.100 = 10.5$؛
و$140 \times 0.116 \approx 16.2$؛ و$175 \times 0.120 = 21.0$؛
و$198 \times 0.120 \approx \qty{23.8}{L/min}$.

\textbf{2.} $23.8/4.5 \approx 5.3$: فالتواتر يسهم بمقدار $198/62
\approx 3.2$، والحجم النفضي بمقدار $120/72 \approx 1.7$، و$3.2 \times
1.7 \approx 5.3$.

\textbf{3.} فوق نحو \qty{225}{W} يبقى الحجم النفضي عند \qty{120}{mL}؛
ولا يرفع الصبيب بعد ذلك إلا التواتر القلبي.

\textbf{4.} $220 - 17 = 203$: و198 عنده ضمن دقة القاعدة.

\textbf{5.} $5/4.5 \approx \qty{1.1}{min}$، أي نحو \qty{67}{s}؛ وعند
\qty{300}{W}، $5/23.8 \approx \qty{0.21}{min}$، أي نحو \qty{13}{s}.

\textbf{6.} $0.5 \times 13 = 6.5$؛ و$1.2 \times 18 = 21.6$؛ و$1.8 \times
24 = 43.2$؛ و$2.3 \times 32 = 73.6$؛ و$2.6 \times 44 \approx
\qty{114}{L/min}$.

\textbf{7.} $114/6.5 \approx 18$: فقد ارتفع الحجم الجاري 5.2 مرات،
والتواتر 3.4 مرات؛ فالحجم يسهم أكثر.

\textbf{8.} $2.6/5.0 = 52\%$ من السعة الحيوية في كل نفس.

\textbf{9.} لا: فللتهوية احتياطي واسع والدم المغادر للرئتين يبقى
مشبعًا. والحد هو الحجم النفضي للقلب، وهو يحدد كم أكسجين توصله كل نبضة.

\textbf{10.} $0.7 \times 6.5 \approx \qty{4.6}{L/min}$ و$0.7 \times
114 \approx \qty{80}{L/min}$.

\textbf{11.} $4.5 \times 0.052 \approx 0.23$؛ و$10.5 \times 0.090
\approx 0.95$؛ و$16.2 \times 0.120 \approx 1.94$؛ و$21.0 \times 0.145
\approx 3.05$؛ و$23.8 \times 0.155 \approx \qty{3.7}{L/min}$.

\textbf{12.} $3700/68 \approx \qty{54}{mL/min}$ لكل كيلوغرام. وقد بلغ
التواتر حده الأقصى والزيادة من 225 إلى \qty{300}{W} صغيرة: فهذا هو
$\dot V\!\mathrm{O_2}$max عنده، وهو تلميذ لائق مدرَّب.

\textbf{13.} $3.7/0.23 \approx 16$: أي عامل القلب 5.3 مضروبًا في عامل
الاستخلاص $155/52 \approx 3.0$.

\textbf{14.} $3.7 \times 20 = \qty{74}{kJ}$ في الدقيقة، أي نحو
\qty{1230}{W}؛ وفي الراحة $0.23 \times 20/60 \approx \qty{77}{W}$؛
وفوق الراحة نحو \qty{1150}{W}، فالمردود $300/1150 \approx
26\%$.

\textbf{15.} خط مستقيم تقريبًا من \num{0.23} إلى \qty{3.05}{L/min} بين
0 و\qty{225}{W}، ثم انحناء: فآخر \qty{75}{W} لا تضيف إلا
\qty{0.65}{L/min}. والسقف نحو \qty{3.7}{L/min}.

\textbf{16.} الاستباق: فتوقع الدماغ للجهد ينقل إلى المراكز القلبية
والغدتين الكظريتين، فيرفع التواتر قبل أي حركة.

\textbf{17.} الأسناخ تؤكسج الدم تأكسيجًا تامًا حتى عند أقصى \omterm{def:g10:heart-lungs-effort:ventilation}{تهوية} وأقصى
صبيب؛ فالرئتان ليستا الحد. والحد هو مقدار الدم الذي يستطيع القلب
إرساله في الدقيقة.

\textbf{18.} $20/23.8 \approx 84\%$. أما المعي والكليتان، وقد كانتا
تأخذان نصف صبيب الراحة، فتتلقيان الآن نحو \qty{1}{L/min}؛ فقد ضاقت
شرايينهما ووجه الجريان إلى العضلات.

\textbf{19.} $198 \times 0.140 \times 0.155 \approx \qty{4.3}{L/min}$،
أي نحو \qty{63}{mL/min} لكل كيلوغرام.

\textbf{20.} تضاعف \omterm{def:g10:heart-lungs-effort:output}{الصبيب القلبي} 5.3 مرات، واستخلاص الأكسجين 3 مرات،
\omterm{def:g10:heart-lungs-effort:ventilation}{والتهوية} 18 مرة؛ وقد غيّر التدريب الحجم النفضي للقلب، ومعه السقف.
\end{solution}
